% !TEX root = ../main.tex
% !TEX program = lualatex
\documentclass[../main.tex]{subfiles}
\IfSubfilesClassLoaded{\externaldocument{\subfix{../build/main}}}{}
% =============================================================================
% File: 25_CNO_Availability_Planning.tex
% =============================================================================
\begin{document}
\IfSubfilesClassLoaded{\chapter{EDO Study Guide}}{}
\section{CNO Availability Planning}

\subsection{Summary}
\begin{description}
  \item[\textbf{Planning balance.}] Availability planning balances flexible requirements with stable early planning to execute within schedule and resources~\autocite{edo-4-2-1-cno-availability-planning-2025}.
  \item[\textbf{AWP sources.}] \ac{awp} work candidates come from \ac{cmp}/\ac{icmp}, authorized alterations, \ac{tycom} repairs (\acp{dfs}/\acp{tso}), pre-availability tests, and \ac{csmp}~\autocite{edo-4-2-1-cno-availability-planning-2025}.
  \item[\textbf{AWP timeline.}] Baseline planning starts at A-36, initial \ac{awp} at A-21, proposed \ac{awp} refinement at mid-planning, post-final \ac{awp} at A-3, and completed \ac{awp} within six months of availability completion~\autocite{edo-4-2-1-cno-availability-planning-2025}.
  \item[\textbf{Work definition.}] Public shipyards use job summaries, component unit phases, and task group instructions; surface availabilities use 4E work item specifications~\autocite{edo-4-2-1-cno-availability-planning-2025}.
  \item[\textbf{Planning guidance.}] \ac{bpmp} and \ac{aim} processes guide baseline project management and industrial planning~\autocite{edo-4-2-1-cno-availability-planning-2025}.
  \item[\textbf{Pre-start goals.}] \ac{lltm}, tooling, early start requirements, critical path, logistics, training, and on-time contract award or definitization must be ready before A-0~\autocite{edo-4-2-1-cno-availability-planning-2025}.
\end{description}

\subsection{Sources of AWP Work Candidates}
\begin{itemize}
  \item \ac{cmp}/\ac{icmp} routine maintenance and inspections~\autocite{edo-4-2-1-cno-availability-planning-2025}.
  \item Authorized alterations (program and fleet)~\autocite{edo-4-2-1-cno-availability-planning-2025}.
  \item \ac{tycom} repairs, including \acp{dfs} and \acp{tso} closures~\autocite{edo-4-2-1-cno-availability-planning-2025}.
  \item Pre-availability tests and inspections (\ac{tsra}, \ac{pots}/\ac{pats})~\autocite{edo-4-2-1-cno-availability-planning-2025}.
  \item \ac{csmp} items~\autocite{edo-4-2-1-cno-availability-planning-2025}.
\end{itemize}

\subsection{Pre-Start Planning Goals}
\begin{itemize}
  \item Order \ac{lltm}, have special tooling on-site, and complete early start requirements~\autocite{edo-4-2-1-cno-availability-planning-2025}.
  \item Define critical path and key events in the integrated schedule~\autocite{edo-4-2-1-cno-availability-planning-2025}.
  \item Resolve logistics and infrastructure needs (power, berthing, galley, force protection, LAN migration)~\autocite{edo-4-2-1-cno-availability-planning-2025}.
  \item Train the crew on shipyard procedures and secure on-time contract award or definitization~\autocite{edo-4-2-1-cno-availability-planning-2025}.
\end{itemize}
The most important planning output is a fully integrated schedule of all availability work~\autocite{edo-4-2-1-cno-availability-planning-2025}.

\subsection{Public Shipyard AWP Development}
\begin{description}
  \item[\textbf{Baseline AWP (A-36).}] Planning requirements from \ac{submepp} (submarines) or \ac{cpaact} (carriers)~\autocite{edo-4-2-1-cno-availability-planning-2025}.
  \item[\textbf{Initial AWP (A-21).}] Consolidates baseline \ac{awp}, authorized alterations, and repairs~\autocite{edo-4-2-1-cno-availability-planning-2025}.
  \item[\textbf{Proposed AWP.}] \ac{tycom} review of pre-availability tests and inspections; \ac{ipm} for submarines and A-12 planning meeting for carriers~\autocite{edo-4-2-1-cno-availability-planning-2025}.
  \item[\textbf{Post-final AWP (A-3).}] \ac{tycom} approval after final planning meeting updates~\autocite{edo-4-2-1-cno-availability-planning-2025}.
  \item[\textbf{Completed AWP.}] Issued within six months of availability completion as requirements gain fidelity over time~\autocite{edo-4-2-1-cno-availability-planning-2025}.
\end{description}

\subsection{RMC AWP Development}
\begin{description}
  \item[\textbf{Baseline AWP (BAWP).}] Developed by \ac{surfmepp} using \ac{cmp}, \ac{tfp}, and ship sheets~\autocite{edo-4-2-1-cno-availability-planning-2025}.
  \item[\textbf{Availability Work Package.}] Refines \ac{bawp} at A-360 with assessment results, \ac{csmp}, and modernization to balance life-cycle maintenance and readiness~\autocite{edo-4-2-1-cno-availability-planning-2025}.
\end{description}

\subsection{Public Shipyard Work Definition}
\begin{description}
  \item[\textbf{Job Summary (JS).}] Groups work into logical evolutions with task lists, material, trades, and man-hour estimates~\autocite{edo-4-2-1-cno-availability-planning-2025}.
  \item[\textbf{Component Unit Phase (CU Phase).}] A specific phase of work for a specific \ac{cu}~\autocite{edo-4-2-1-cno-availability-planning-2025}.
  \item[\textbf{Task Group Instruction (TGI).}] Detailed instructions for tasks within a \ac{cu} phase~\autocite{edo-4-2-1-cno-availability-planning-2025}.
  \item[\textbf{Component Unit (CU).}] The physical unit of work control~\autocite{edo-4-2-1-cno-availability-planning-2025}.
\end{description}
\begin{description}
  \item[\textbf{CU phases.}] A (repair onboard), C (calibrate/align), D (open/inspect), E (restore prior to test), F (fabricate), H (repair off-ship), I (inspect), K (clean/prepare), M (maintain), P (paint/preserve), R (reinstall/reassemble), S (stage/prepare), T (test), U (unship/remove)~\autocite{edo-4-2-1-cno-availability-planning-2025}.
\end{description}

\subsection{Integrated Schedule (Public Shipyard)}
\begin{enumerate}
  \item Define the authorized \ac{awp} and establish key events and milestones~\autocite{edo-4-2-1-cno-availability-planning-2025}.
  \item Develop job summaries and sequence \ac{cu} phases~\autocite{edo-4-2-1-cno-availability-planning-2025}.
  \item Build the network diagram (PERT), integrate Ship's Force/contractor/\ac{ait} work, and estimate resources, time, and cost~\autocite{edo-4-2-1-cno-availability-planning-2025}.
  \item Determine the critical path and produce the integrated schedule~\autocite{edo-4-2-1-cno-availability-planning-2025}.
\end{enumerate}

\subsection{Critical and Controlling Paths}
\begin{description}
  \item[\textbf{Critical path.}] Longest duration sequence of work with the least float that sets minimum project length~\autocite{edo-4-2-1-cno-availability-planning-2025}.
  \item[\textbf{Controlling path.}] Near-critical activities with slightly more float that can still drive key events due to scope, material, or complexity~\autocite{edo-4-2-1-cno-availability-planning-2025}.
\end{description}

\subsection{RMC Work Definition and Schedule}
\begin{description}
  \item[\textbf{4E work item specifications.}] Contractual statements of work defining repair or alteration requirements in standard format for Master Ship Repair Agreement or Agreement for Boat Repair~\autocite{edo-4-2-1-cno-availability-planning-2025}.
  \item[\textbf{Integrated schedule.}] Defines \ac{awp}, phases, key events, milestones, network sequence, and critical path; \ac{lma} generates the schedule, contracted through \ac{navsea} Standard Item 009-60, with \ac{rmc} review and acceptance~\autocite{edo-4-2-1-cno-availability-planning-2025}.
\end{description}

\subsection{Platform Differences}
\begin{description}
  \item[\textbf{Baseline planning owners.}] \ac{submepp} (submarines), \ac{cpaact} (carriers), and \ac{surfmepp} (surface) maintain the baseline \ac{awp}~\autocite{edo-4-2-1-cno-availability-planning-2025}.
  \item[\textbf{Planning organizations.}] Submarines and carriers rely on public/private shipyard planning products; surface ships use contractor or third-party \ac{sdaes} planning with 4E specifications~\autocite{edo-4-2-1-cno-availability-planning-2025}.
  \item[\textbf{Common sources.}] All platforms use the same work candidate sources, but planning timelines differ across platforms~\autocite{edo-4-2-1-cno-availability-planning-2025}.
\end{description}

\subsection{Planning Team Roles}
\begin{description}
  \item[\textbf{BSPO.}] Customer interface and funding authority; processes new work and funding~\autocite{edo-4-2-1-cno-availability-planning-2025}.
  \item[\textbf{Project superintendent.}] Leads cost, schedule, and performance; oversees project planning and sequencing~\autocite{edo-4-2-1-cno-availability-planning-2025}.
  \item[\textbf{PEPM.}] Principal planner who directs \ac{js} boundaries, planning products, and schedule development~\autocite{edo-4-2-1-cno-availability-planning-2025}.
  \item[\textbf{Project engineers/test engineers.}] Support planning, sequencing, and technical oversight~\autocite{edo-4-2-1-cno-availability-planning-2025}.
  \item[\textbf{Port engineer.}] \ac{tycom} representative who screens work to the proper maintenance organization~\autocite{edo-4-2-1-cno-availability-planning-2025}.
  \item[\textbf{RMC project manager.}] Leads the RMC project team, ensures funding for planning, and screens work to contractor execution~\autocite{edo-4-2-1-cno-availability-planning-2025}.
  \item[\textbf{LMA project manager.}] Contractor \ac{lma} point of contact who leads integrated schedule development and coordinates with the \ac{rmc}~\autocite{edo-4-2-1-cno-availability-planning-2025}.
  \item[\textbf{SMMO.}] Ship's Force point of contact for maintenance; assists with screening and 2K development~\autocite{edo-4-2-1-cno-availability-planning-2025}.
\end{description}

\subsection{RMC Contracting Strategy (MACMO)}
\begin{description}
  \item[\textbf{MAC-MO.}] Firm-fixed-price surface ship repair contracts that separate planning from execution (\ac{sdaes} or organic planning), maximize competition, and balance risk between Navy and industry~\autocite{edo-4-2-1-cno-availability-planning-2025}.
  \item[\textbf{Award approach.}] \ac{cno} availabilities of 10+ months compete coast-wide; under 10 months award in homeport; continuous maintenance uses homeport competition; emergent work uses existing contractor, competition if time allows, or rotational \ac{idiq} awards if immediate~\autocite{edo-4-2-1-cno-availability-planning-2025}.
\end{description}

\subsection{Quick Review}
\begin{itemize}
  \item Five \ac{awp} sources: \ac{cmp}/\ac{icmp}, authorized alterations, \ac{tycom} repairs, pre-availability tests, and \ac{csmp}~\autocite{edo-4-2-1-cno-availability-planning-2025}.
  \item Baseline \ac{awp} producers: \ac{submepp} (subs), \ac{cpaact} (carriers), \ac{surfmepp} (surface)~\autocite{edo-4-2-1-cno-availability-planning-2025}.
  \item The most important planning product is the integrated schedule~\autocite{edo-4-2-1-cno-availability-planning-2025}.
  \item \ac{macmo} uses fixed-price competition to drive cost down and improve accountability~\autocite{edo-4-2-1-cno-availability-planning-2025}.
\end{itemize}

\IfSubfilesClassLoaded{
  \RenewDocumentCommand{\entryname}{}{\textbf{\color{Modern} Acronym}}
  \RenewDocumentCommand{\descriptionname}{}{\textbf{\color{Modern} Definition}}
  \printnoidxglossary[
    type=\acronymtype,
    title=Acronyms,
    style=long-booktabs
  ]}{}
\subsection{Coursebook cue: Team One comparison and planning products}

Coursebook 4.2.1 compares submarine, carrier, and surface planning to show that
the planning logic is common even when the owning organizations and products
differ
\autocite{edo-4-2-1-cno-availability-planning-2025}.

\begin{longtblr}[
  caption = {Availability planning comparison.},
  label = {tab:availability-planning-comparison},
]{
  width = \linewidth,
  colspec = {X[0.65,l] X[0.95,l] X[1.2,l]},
  rowhead = 1,
  hlines,
  vlines,
}
\textbf{Community} & \textbf{Baseline AWP owner} & \textbf{Planning cue} \\
Submarine & SUBMEPP & Public or private shipyard plans and sequences work using job summaries and task group instructions. \\
Carrier & Carrier Planning Activity & Carrier Team One focuses on aircraft-carrier maintenance performance and uses similar work-package logic. \\
Surface & SURFMEPP & Surface Team One focuses on surface-ship service-life performance; planning may use contractor or third-party Specification Development and Availability Execution Support with 4E work-item specifications. \\
\end{longtblr}

\begin{longtblr}[
  caption = {Availability planning products.},
  label = {tab:availability-planning-products},
]{
  width = \linewidth,
  colspec = {X[0.75,l] X[1.75,l]},
  rowhead = 1,
  hlines,
  vlines,
}
\textbf{Product} & \textbf{Use} \\
Availability Work Package & Authoritative list of planned work for the availability. \\
Job Summary & Planning record containing the information needed to plan and schedule a work evolution. \\
Control Unit phase & Phase of work performed on a specific control unit. \\
Task Group Instruction & Detailed work instructions and technical information for a control-unit phase. \\
4E work-item specification & Surface-ship statement of work for the planned repair or modernization item. \\
Integrated schedule & Networked schedule that links authorized work, key events, critical path, and execution constraints. \\
\end{longtblr}

\end{document}
